\scanpage{35}
(i)
\[
\begin{aligned}
 F&=\deleted{\unsure{X^p+Y^p-a}}aX^p+bY^p-1\\
 k[X,Y]/(F)&=A\\[-2pt]
 &\unsure{\wr\!\mid|\simeq}\\[-2pt]
 &k[Z]/(Z^p-a)[X]
\end{aligned}
\]
\margin{faire $Z=X+Y$ donc $F=Z^p-a$}
\margin{$aX^p+bY^p$}

\deleted{\unsure{transcendans}} $A$ est un anneau régulier, \inserted{non} $k$ est alg.\textsuperscript{t} fermé dans son corps des fractions

$A\otimes_k\bar k$ n'est pas intègre, il est isomorphe à $k[Z][T]/(T^p)$

\noindent\rule{.28\linewidth}{.3pt}

(ii)
\[
\left.
\begin{aligned}
 F&=Y^2-X^p+a\\
 k[X,Y]/(F)&=A
\end{aligned}
\right|\,p\ne2
\]
$A$ est un anneau régulier, \deleted{\unsure{$k$ est alg.\textsuperscript{t}}} son corps des fractions est extension séparable de $k$

$A\otimes_k\bar k$ n'est pas régulier, car le point $Y=0$, $X=a^{1/p}$ est singulier sur la courbe $F=0$

En effet, la courbe devient $F=Y^2-(X-\lambda)^p=0$ \quad$(\lambda=a^{1/p})$. Faisons $X-\lambda=Z$, on a donc $k[Y,Z]/(Y^2-Z^p)$. Or $Y^2-Z^p\in{\unsure{m|\mathfrak m}}^{2}$, où $m$ est l'idéal maximal de l'origine, donc on ne peut avoir point simple à l'origine. Cependant, pour $p\ne\unsure{2|1}$, $Y^2-Z^p$ est un polynôme irréductible. Cela prouve en $=$ temps que $K/k$ est une extension séparable. Montrons que $A$ est régulier.

\deleted{Dans la courbe obtenue par extension à $\bar k$, on voit que le seul point singulier est} \unsure{Car|le} .

$\Omega^1_k(A)\simeq A^2/A\bigl(\frac{\partial F}{\partial X},\frac{\partial F}{\partial Y}\bigr)\simeq A+A/2Y$. Les seuls points où $\Omega^1_k(A)$ soit non libre est aux idéaux premiers $\supset Y$, i.e. les idéaux premiers de $k[X,Y]$ \unsure{cont\textsuperscript{t}} $(Y,\deleted{\unsure{a}}Y^2-X^p+a)=(Y,X^p-a)$ ; or $k[X,Y]/(Y,X^p-a)\simeq k[X]/(X^{p}-a)$

Donc l'idéal $m'=(Y,X^p-a)$ est déjà maximal, c'est donc le

% archive sheet number: 34

